% !TeX root = ../main.tex
\section{Introduction}

Software architecture describes a system's intended components, responsibilities, and permitted relationships. Implementation can diverge from this intent. Mapping work reports structural violations and effects on software quality and evolution; developer evidence identifies continuing reliance on documentation, clear responsibilities, and shared practices \cite{li2022erosion,anthony2024drifting}. A local change can bypass a gateway, couple an interface directly to a database, or add infrastructure without an architecture decision while compilation and functional tests pass.

Without executable checks, documentation can leave a gap between intended and implemented structure. ArchSync connects architecture intent and implementation through deterministic, evidence-backed checks. Diagrams communicate intent but are difficult to use as automated gates; source-derived graphs expose relationships but do not specify which are allowed, forbidden, required, or subject to review.

We present ArchSync, an architecture-conformance workflow based on two graphs. The expected graph is derived from a versioned \texttt{architecture.yaml} document containing components, relationships, and deterministic rules. The observed graph is inferred from source code by stack-specific analyzers. ArchSync compares the graphs and produces one of three classifications: \emph{no-impact} when implementation conforms without topology changes, \emph{violation} when one or more rules are broken, and \emph{evolution} when a non-forbidden topology change requires human review. Findings retain both the model location that establishes the expectation and the source location supporting the observation.

For pull requests, the workflow compares findings at the Git merge base with findings at the proposed head. Only findings introduced by the change determine its gate decision: violations block, non-forbidden topology changes request review, and changes without new drift pass. The observed graph therefore follows implementation changes, while the expected graph remains the approved design baseline. Our proposed governance workflow requires an explicit model change and architecture-owner review to accept evolution; the evaluated command reports the need for review but does not manage approval or rejection records.

The prototype deliberately focuses on TypeScript and Node.js. It recognizes HTTP calls through \texttt{fetch}, PostgreSQL access through \texttt{pg}, Redis operations, and AMQP publishing and consumption. This is an empirical boundary rather than a claim that architecture drift is TypeScript-specific: the expected graph, rule engine, finding contract, and metrics are language-neutral, while each additional language requires a separately implemented and evaluated source-to-graph analyzer. The bounded TypeScript scope permits deterministic contracts and reproducible measurement without probabilistic inference. Large-language-model explanations and automatic repairs are excluded from the decision path.

This paper makes the following contributions:
\begin{itemize}
  \item an executable-contract prototype integrating typed service relationships, four conformance rules, and source/model evidence;
  \item a relation-level Git-diff gate that separates introduced violations, reviewable evolution, and historical findings using cached component-scoped observation;
  \item a reproducible controlled verification artifact comprising 20 patches and 40 detector signals, with explicit units, provenance, and limits of inference.
\end{itemize}
These are bounded prototype contributions, not claims that expected-versus-observed comparison, dependency rules, or deterministic analysis are new concepts.
